% GENERATED FILE. Edit canonical lessons, then run npm run generate:books.
% canonical-source-hash: fnv1a64:92caf37a7b606431
% canonical-lessons: ML-C430-alimati, ML-C430-kaikkuli, ML-C430-cusanam-ceyyuka, ML-C430-sutaryamaya, ML-C430-anvesikkuka

\chapter{Corruption}
\label{ch:ml-corruption}
\hlchaptermodality{430}

\begin{chapteropening}
\textbf{By the end of this chapter:} \emph{I can say corruption, a bribe, to exploit, transparent and to investigate.}

Complete the last lesson of chapter 430: I can say corruption, a bribe, to exploit, transparent and to investigate.
\end{chapteropening}

\section[\texorpdfstring{\ml{അഴിമതി} (aḻimati)}{അഴിമതി (aḻimati)}]{\ml{അഴിമതി} (aḻimati) — corruption}

\label{lesson:ML-C430-alimati}



\begin{quote}
Before the new one: say the Malayalam for a benefit, then the Malayalam for to distribute.
\end{quote}

\subsection*{You'll want to know: \ml{അഴിമതി}}
\textbf{\ml{അഴിമതി}} — \emph{aḻimati} — \textquotedblleft{}corruption\textquotedblright{}.

\begin{quote}
\ml{ആ} \ml{ഓഫീസിൽ} \ml{വലിയ} \ml{അഴിമതി} \ml{നടന്നു}.
\end{quote}

\begin{quote}
\emph{ā ōphīsil valiya aḻimati naṭannu.}
\end{quote}

\begin{quote}
There was large-scale corruption in that office.
\end{quote}

Corruption is \ml{നടക്കുക}, carried on.

\begin{cousinweb}
\textbf{\ml{അഴിമതി}} (\emph{aḻimati}) is a Malayalam word whose origin is not settled; DEDR has no entry for it.
\end{cousinweb}

\begin{grammarlens}[title={What you've built}]
One more word for giving reasons and opinions.
\end{grammarlens}

\subsection*{Your turn}
\begin{itemize}
\raggedright
  \item \emph{Say it:} \emph{aḻimati}
  \item \emph{Say it:} \emph{aḻimati}, once more
  \item \emph{Say it:} \emph{ā ōphīsil valiya aḻimati naṭannu.} — There was large-scale corruption in that office.
  \item \emph{Recall:} say the Malayalam for a benefit, then the Malayalam for to distribute, then say \emph{aḻimati} again
\end{itemize}

\begin{tcolorbox}[breakable,colback=teal!4,colframe=teal!35!black,title={Before you move on}]
What does \ml{അഴിമതി} mean? (\textquotedblleft{}Corruption\textquotedblright{}.) Say it once more.
\end{tcolorbox}
\section[\texorpdfstring{\ml{കൈക്കൂലി} (kaikkūli)}{കൈക്കൂലി (kaikkūli)}]{\ml{കൈക്കൂലി} (kaikkūli) — a bribe}

\label{lesson:ML-C430-kaikkuli}



\begin{quote}
Before the new one: say the Malayalam for to distribute, then the Malayalam for corruption.
\end{quote}

\subsection*{You'll want to know: \ml{കൈക്കൂലി}}
\textbf{\ml{കൈക്കൂലി}} — \emph{kaikkūli} — \textquotedblleft{}a bribe\textquotedblright{}.

\begin{quote}
\ml{ആ} \ml{ഉദ്യോഗസ്ഥൻ} \ml{കൈക്കൂലി} \ml{വാങ്ങി}.
\end{quote}

\begin{quote}
\emph{ā udyōgasthan kaikkūli vāṅṅi.}
\end{quote}

\begin{quote}
That official took a bribe.
\end{quote}

To take a bribe is \ml{കൈക്കൂലി} \ml{വാങ്ങുക}; to give one, \ml{കൈക്കൂലി} \ml{കൊടുക്കുക}.

\begin{cousinweb}
\textbf{\ml{കൈക്കൂലി}} (\emph{kaikkūli}) is a native Malayalam compound: \ml{കൈ}, \textquotedblleft{}hand\textquotedblright{} (DEDR 2023), and \ml{കൂലി}, \textquotedblleft{}wages\textquotedblright{} (DEDR 1905). DEDR 1905 lists \emph{kai-kkūli} itself with the sense \textquotedblleft{}bribe\textquotedblright{}: pay slipped into the hand.
\end{cousinweb}

\begin{grammarlens}[title={What you've built}]
One more word for giving reasons and opinions.
\end{grammarlens}

\subsection*{Your turn}
\begin{itemize}
\raggedright
  \item \emph{Say it:} \emph{kaikkūli}
  \item \emph{Say it:} \emph{kaikkūli}, once more
  \item \emph{Say it:} \emph{ā udyōgasthan kaikkūli vāṅṅi.} — That official took a bribe.
  \item \emph{Recall:} say the Malayalam for to distribute, then the Malayalam for corruption, then say \emph{kaikkūli} again
\end{itemize}

\begin{tcolorbox}[breakable,colback=teal!4,colframe=teal!35!black,title={Before you move on}]
What does \ml{കൈക്കൂലി} mean? (\textquotedblleft{}A bribe\textquotedblright{}.) Say it once more.
\end{tcolorbox}
\section[\texorpdfstring{\ml{ചൂഷണം} \ml{ചെയ്യുക} (cūṣaṇaṁ ceyyuka)}{ചൂഷണം ചെയ്യുക (cūṣaṇaṁ ceyyuka)}]{\ml{ചൂഷണം} \ml{ചെയ്യുക} (cūṣaṇaṁ ceyyuka) — to exploit}

\label{lesson:ML-C430-cusanam-ceyyuka}



\begin{quote}
Before the new one: say the Malayalam for corruption, then the Malayalam for a bribe.
\end{quote}

\subsection*{You'll want to know: \ml{ചൂഷണം} \ml{ചെയ്യുക}}
\textbf{\ml{ചൂഷണം} \ml{ചെയ്യുക}} — \emph{cūṣaṇaṁ ceyyuka} — \textquotedblleft{}to exploit\textquotedblright{}.

\begin{quote}
\ml{ആ} \ml{കമ്പനി} \ml{തൊഴിലാളികളെ} \ml{ചൂഷണം} \ml{ചെയ്തു}.
\end{quote}

\begin{quote}
\emph{ā kampani toḻilāḷikaḷe cūṣaṇaṁ ceytu.}
\end{quote}

\begin{quote}
That company exploited its workers.
\end{quote}

Past: \ml{ചൂഷണം} \ml{ചെയ്തു}, \emph{cūṣaṇaṁ ceytu}.

The noun \ml{ചൂഷണം} is exploitation.

\begin{cousinweb}
\textbf{\ml{ചൂഷണം} \ml{ചെയ്യുക}} (\emph{cūṣaṇaṁ ceyyuka}) is a phrase verb built on Sanskrit \emph{cūṣaṇa}, \textquotedblleft{}sucking\textquotedblright{} (CDIAL 4897), with \ml{ചെയ്യുക}, \textquotedblleft{}to do\textquotedblright{}. To exploit is to suck a person dry.
\end{cousinweb}

\begin{grammarlens}[title={What you've built}]
One more word for giving reasons and opinions.
\end{grammarlens}

\subsection*{Your turn}
\begin{itemize}
\raggedright
  \item \emph{Say it:} \emph{cūṣaṇaṁ ceyyuka}
  \item \emph{Say it:} \emph{cūṣaṇaṁ ceyyuka}, once more
  \item \emph{Say it:} \emph{ā kampani toḻilāḷikaḷe cūṣaṇaṁ ceytu.} — That company exploited its workers.
  \item \emph{Recall:} say the Malayalam for corruption, then the Malayalam for a bribe, then say \emph{cūṣaṇaṁ ceyyuka} again
\end{itemize}

\begin{tcolorbox}[breakable,colback=teal!4,colframe=teal!35!black,title={Before you move on}]
What does \ml{ചൂഷണം} \ml{ചെയ്യുക} mean? (\textquotedblleft{}To exploit\textquotedblright{}.) Say it once more.
\end{tcolorbox}
\section[\texorpdfstring{\ml{സുതാര്യമായ} (sutāryamāya)}{സുതാര്യമായ (sutāryamāya)}]{\ml{സുതാര്യമായ} (sutāryamāya) — transparent}

\label{lesson:ML-C430-sutaryamaya}



\begin{quote}
Before the new one: say the Malayalam for a bribe, then the Malayalam for to exploit.
\end{quote}

\subsection*{You'll want to know: \ml{സുതാര്യമായ}}
\textbf{\ml{സുതാര്യമായ}} — \emph{sutāryamāya} — \textquotedblleft{}transparent\textquotedblright{}.

\begin{quote}
\ml{ഭരണം} \ml{സുതാര്യമായിരിക്കണം}.
\end{quote}

\begin{quote}
\emph{bharaṇaṁ sutāryamāyirikkaṇaṁ.}
\end{quote}

\begin{quote}
Administration must be transparent.
\end{quote}

It is said of glass and of honest government alike.

\begin{cousinweb}
\textbf{\ml{സുതാര്യമായ}} (\emph{sutāryamāya}) is a modern learned Malayalam word whose origin is not settled.
\end{cousinweb}

\begin{grammarlens}[title={What you've built}]
One more word for giving reasons and opinions.
\end{grammarlens}

\subsection*{Your turn}
\begin{itemize}
\raggedright
  \item \emph{Say it:} \emph{sutāryamāya}
  \item \emph{Say it:} \emph{sutāryamāya}, once more
  \item \emph{Say it:} \emph{bharaṇaṁ sutāryamāyirikkaṇaṁ.} — Administration must be transparent.
  \item \emph{Recall:} say the Malayalam for a bribe, then the Malayalam for to exploit, then say \emph{sutāryamāya} again
\end{itemize}

\begin{tcolorbox}[breakable,colback=teal!4,colframe=teal!35!black,title={Before you move on}]
What does \ml{സുതാര്യമായ} mean? (\textquotedblleft{}Transparent\textquotedblright{}.) Say it once more.
\end{tcolorbox}
\section[\texorpdfstring{\ml{അന്വേഷിക്കുക} (anvēṣikkuka)}{അന്വേഷിക്കുക (anvēṣikkuka)}]{\ml{അന്വേഷിക്കുക} (anvēṣikkuka) — to investigate; to inquire}

\label{lesson:ML-C430-anvesikkuka}



\begin{quote}
Before the new one: say the Malayalam for to exploit, then the Malayalam for transparent.
\end{quote}

\subsection*{You'll want to know: \ml{അന്വേഷിക്കുക}}
\textbf{\ml{അന്വേഷിക്കുക}} — \emph{anvēṣikkuka} — \textquotedblleft{}to investigate; to inquire\textquotedblright{}.

\begin{quote}
\ml{പോലീസ്} \ml{ആ} \ml{സംഭവം} \ml{അന്വേഷിക്കുന്നു}.
\end{quote}

\begin{quote}
\emph{pōlīsŭ ā sambhavaṁ anvēṣikkunnu.}
\end{quote}

\begin{quote}
The police are investigating that incident.
\end{quote}

Past: \ml{അന്വേഷിച്ചു}, \emph{anvēṣiccu}.

Also to look for, to ask after: \ml{അവൻ} \ml{നിന്നെ} \ml{അന്വേഷിച്ചു}, he asked after you.

\begin{cousinweb}
\textbf{\ml{അന്വേഷിക്കുക}} (\emph{anvēṣikkuka}) is built from Sanskrit \emph{anv-eṣaṇa}, \textquotedblleft{}seeking after, searching\textquotedblright{} (Monier-Williams), with the Malayalam verb ending -\ml{ിക്കുക}.
\end{cousinweb}

\begin{grammarlens}[title={What you've built}]
One more word for giving reasons and opinions.
\end{grammarlens}

\subsection*{Your turn}
\begin{itemize}
\raggedright
  \item \emph{Say it:} \emph{anvēṣikkuka}
  \item \emph{Say it:} \emph{anvēṣikkuka}, once more
  \item \emph{Say it:} \emph{pōlīsŭ ā sambhavaṁ anvēṣikkunnu.} — The police are investigating that incident.
  \item \emph{Recall:} say the Malayalam for to exploit, then the Malayalam for transparent, then say \emph{anvēṣikkuka} again
\end{itemize}

\begin{tcolorbox}[breakable,colback=teal!4,colframe=teal!35!black,title={Before you move on}]
What does \ml{അന്വേഷിക്കുക} mean? (\textquotedblleft{}To investigate; to inquire\textquotedblright{}.) Say it once more.
\end{tcolorbox}
